\documentclass[11pt]{article}
\usepackage[margin=1in]{geometry}
\usepackage{amsmath,amssymb,amsthm}
\usepackage{graphicx}
\usepackage{microtype}
\usepackage[hidelinks]{hyperref}

\newtheorem{theorem}{Theorem}
\newtheorem{lemma}{Lemma}
\newcommand{\Hh}{\mathcal H}
\newcommand{\Qlin}{Q}
\newcommand{\Lin}{\operatorname{Lin}}
\newcommand{\dist}{\operatorname{dist}}
\newcommand{\spanop}{\operatorname{span}}

\title{Macaev Liftability Is Equivalent to\\Logarithmic Local Linearization}
\author{Automated mathematics run \texttt{fa\_banach\_001}, lane 9}
\date{11 August 2026}

\begin{document}
\maketitle

\noindent\fbox{\parbox{0.94\textwidth}{\textbf{Status.} Candidate partial
result, likely valid. The universal problem is reduced exactly to a sharp
finite-dimensional estimate, but that estimate is not decided here.}}

\section{Source question}

Cabello Sanchez and Garcia ask on page 5 of \cite{CG} whether every operator
in the Macaev ideal is liftable through every twisted Hilbert extension.
Their Macaev ideal consists of compact operators $T$ such that
\[
  \sum_{j\geq1}\frac{s_j(T)}j<\infty.
\]

\begin{figure}[h]
  \centering
  \includegraphics[width=0.82\textwidth]{figures/open_problem_crop.png}
  \caption{The source question in arXiv:2112.03852, page 5.}
\end{figure}

For a homogeneous quasilinear map $\Phi:E\to F$, let $\Qlin(\Phi)$ be the
least constant in
\[
 \|\Phi(x+y)-\Phi(x)-\Phi(y)\|
 \leq \Qlin(\Phi)(\|x\|+\|y\|),
\]
and put
\[
 \dist(\Phi,\Lin)
 =\inf_{L\text{ linear}}\sup_{x\ne0}
   \frac{\|\Phi(x)-Lx\|}{\|x\|}.
\]
Fixing a separable infinite-dimensional Hilbert space $\Hh$, define
\[
 \Lambda(m)=\sup\bigl\{\dist(\Phi,\Lin):
 \Phi:\ell_2^m\to\Hh,\ \Qlin(\Phi)\leq1\bigr\}.
\]
This number is finite. Indeed, subtract the linear map agreeing with $\Phi$
on an orthonormal basis and combine the coordinates in a balanced binary
tree. At each level Cauchy--Schwarz bounds the sum of the norms of the
partial coordinate blocks. Summing the resulting geometric series gives,
for example, $\Lambda(m)\leq4\sqrt m$.

\section{Idea of the proof}

The positive direction is Abel summation. A bound
$\Lambda(m)=O(\log(m+1))$ gives logarithmically bounded liftings of rank-$m$
projections. The spectral decomposition
\[
 D=\sum_{m\geq1}(s_m-s_{m+1})P_m
\]
then becomes a convergent series of liftings precisely because the Macaev
condition is the discrete integral of those logarithmic costs.

The converse uses the same mechanism in reverse. If the logarithmic bound
fails, select increasingly bad finite-dimensional quasilinear maps and place
them on orthogonal blocks. Scalar coefficients are chosen so that their
Macaev contributions are summable while their block linearization
obstructions tend to infinity. The resulting block-scalar Macaev operator
cannot lift through the completed block twist.

\section{The endpoint dichotomy}

\begin{theorem}\label{thm:main}
The following are equivalent.
\begin{enumerate}
\item Every operator in the Macaev ideal on a separable Hilbert space is
liftable through every twisted Hilbert extension.
\item There is a universal constant $C$ such that
\[
  \Lambda(m)\leq C\log(m+1)\qquad(m\geq1).
\]
\end{enumerate}
More precisely, if (2) fails, a block direct-sum twist and a block-scalar
Macaev operator give a counterexample to (1).
\end{theorem}

\begin{proof}
Assume (2). By the ideal property of liftable operators \cite[Lemma 1]{CG},
it is enough to treat a positive diagonal compact operator $D$ with decreasing
diagonal $(s_m)$. Let $P_m$ be the projection onto its first $m$ eigenvectors
and set $d_m=s_m-s_{m+1}$. Thus
\[
 D=\sum_{m\geq1}d_mP_m.
\]

Fix a twisted Hilbert extension and represent it by a quasilinear map
$\Phi:\Hh\to\Hh$. Restricting $\Phi$ to $P_m\Hh$ and using (2), choose a
linear map $L_m:P_m\Hh\to\Hh$ such that
\[
 \|\Phi(P_mx)-L_mP_mx\|
 \leq C\Qlin(\Phi)\log(m+1)\|x\|.
\]
Hence $x\mapsto(L_mP_mx,P_mx)$ is a lifting of $P_m$ whose norm, in the
standard twisted-sum quasinorm, is at most
$1+C\Qlin(\Phi)\log(m+1)$.

Summation by parts gives
\begin{align*}
 \sum_{m\geq1}d_m\log(m+1)
 &=s_1\log2+
   \sum_{m\geq2}s_m\bigl(\log(m+1)-\log m\bigr)\\
 &\leq s_1\log2+\sum_{m\geq2}\frac{s_m}{m}<\infty.
\end{align*}
A twisted Hilbert space has an equivalent Banach norm. Therefore the weighted
series of the displayed liftings converges to a bounded linear lifting of
$D$. This proves (1).

We prove the converse contrapositively. Suppose that
$\Lambda(m)/\log(m+1)$ is unbounded. Choose strictly increasing integers
$m_k$ and maps
\[
 \Phi_k:\ell_2^{m_k}\to Y_k,\qquad \Qlin(\Phi_k)\leq1,
\]
where the $Y_k$ are Hilbert spaces, so that
\[
 \delta_k:=\dist(\Phi_k,\Lin)
 \geq k^4\log(m_k+1).
\]
On the algebraic Hilbertian direct sum define
\[
 \Omega((x_k))=(\Phi_k(x_k)).
\]
Minkowski's inequality shows that $\Qlin(\Omega)\leq1$. Completing the
associated twisted sum gives an exact sequence whose quotient is
$X=(\bigoplus_k\ell_2^{m_k})_2$ and whose kernel is
$Y=(\bigoplus_kY_k)_2$.

Set
\[
 a_k=\frac1{k^2\log(m_k+1)},\qquad
 D=\bigoplus_{k\geq1}a_k I_{\ell_2^{m_k}}.
\]
The numbers $a_k$ decrease to zero, so $D$ is compact. Put
$N_k=\sum_{i\leq k}m_i$. Since the singular value $a_k$ occurs $m_k$
times, the contribution of the $k$-th block, for $k\geq2$, satisfies
\[
 \sum_{j=N_{k-1}+1}^{N_k}\frac{s_j(D)}j
 =a_k\sum_{j=N_{k-1}+1}^{N_k}\frac1j
 \leq a_k\log\frac{N_k}{N_{k-1}}
 \leq a_k\log(m_k+1)=\frac1{k^2}.
\]
The first block contributes only a finite amount, and therefore $D$ belongs
to the Macaev ideal.

If $D$ had a bounded lifting through the completed block twist, restriction
to the $k$-th quotient block and projection onto the $k$-th kernel block
would give a linear $L_k:\ell_2^{m_k}\to Y_k$ with
\[
  \sup_k\|a_k\Phi_k-L_k\|<\infty.
\]
By homogeneity this forces $\sup_k a_k\delta_k<\infty$. But
\[
  a_k\delta_k\geq k^2,
\]
a contradiction. Thus $D$ is a Macaev operator which is not liftable, so
(1) fails.
\end{proof}

\section{A sharp finite-range theorem}

The full logarithmic estimate remains undecided when the target Hilbert space
is unrestricted.  There is, however, a sharp theorem when both dimensions are
controlled.  Let
\[
 \Lambda_{\rm c}(m,r)=\sup\bigl\{\dist(\Phi,\Lin):
 \Phi:\ell_2^m\to\ell_2^r\text{ continuous and homogeneous},\
 \Qlin(\Phi)\leq1\bigr\}.
\]

\begin{theorem}\label{thm:finite-range}
There is a universal constant $C$ such that
\[
 \Lambda_{\rm c}(m,r)\leq C\bigl(1+\log(m+r)\bigr).
\]
This order is sharp when $m$ and $r$ are comparable.  Consequently, fix a
continuous quasilinear map $\Phi:\Hh\to\Hh$ and nested finite-dimensional
subspaces $E_m\subset\Hh$, $\dim E_m=m$.  If, for some fixed $a<\infty$,
\[
 \dim\spanop\Phi(E_m)\leq(m+1)^a\qquad(m\geq1),
\]
then every positive diagonal Macaev operator whose first $m$ eigenvectors span
$E_m$ lifts through the twist generated by $\Phi$.
\end{theorem}

\begin{proof}
Put $n=\max(m,r)$, embed domain and range isometrically in $\ell_2^n$, and let
$P:\ell_2^n\to\ell_2^m$ be the orthogonal projection.  The map
\[
 \widetilde\Phi(x)=J\Phi(Px):\ell_2^n\longrightarrow\ell_2^n
\]
is continuous, homogeneous, and has quasilinearity constant at most one.
Kalton's logarithmic linearization theorem
\cite[Theorem~2.2]{Kalton2002} supplies a linear $A:\ell_2^n\to\ell_2^n$
such that
\[
 \|\widetilde\Phi(x)-Ax\|
 \leq C(1+\log n)\|x\|.
\]
Restricting $A$ to the embedded $\ell_2^m$ and orthogonally projecting its
values to the embedded $\ell_2^r$ proves the upper bound.  Kalton's example
\cite[Proposition~2.1]{Kalton2002} gives the matching logarithmic lower order
when the dimensions are comparable.

For the consequence, restrict $\Phi$ to $E_m$ and identify the Hilbert space
$\spanop\Phi(E_m)$ with $\ell_2^{r_m}$, where $r_m\leq(m+1)^a$.  The first
part yields linear maps $L_m:E_m\to\Hh$ satisfying
\[
 \|\Phi(x)-L_mx\|
 \leq C_a\Qlin(\Phi)\log(m+1)\|x\|\qquad(x\in E_m).
\]
The Abel-summation proof in Theorem~\ref{thm:main} now applies verbatim.
\end{proof}

\section{Verification and limitation}

The proof uses only two-sided estimates: Abel summation for the positive
direction and restriction to individual blocks for the negative direction.
The dense block map is used only to define its completed twisted sum, exactly
as in standard block constructions of twisted Hilbert spaces. No uniform
boundedness of the individual maps is assumed or needed.

The endpoint dichotomy identifies the remaining issue sharply, but does not
decide it.  Theorem~\ref{thm:finite-range} shows that target dimension is the
only unresolved quantitative parameter in the continuous finite-dimensional
problem: Kalton's theorem gives $O(\log(m+r))$, whereas the universal Macaev
question requires $O(\log m)$ with no restriction on $r$, including an
infinite-dimensional target.  A dualization followed by Maurey's projection
theorem does not remove this parameter: estimating the required type constant
is quantitatively equivalent to bounding the original splitting constant.
The available Schatten-class argument in \cite{CG} supplies a bound depending
on every fixed exponent $p<\infty$; its cited splitting theorem is
nonquantitative in $p$. Consequently it does not currently yield the linear
growth in $p$ that would imply (2). Human review should concentrate on the
completed block construction and on whether an existing local theory theorem
already decides the asymptotic behavior of $\Lambda(m)$.

\begin{thebibliography}{9}
\bibitem{CG}
F. Cabello Sanchez and R. Garcia,
\emph{The Twisted Hilbert Space Ideals},
Integral Equations Operator Theory 94 (2022), article 21;
arXiv:2112.03852.

\bibitem{Kalton2002}
N. J. Kalton,
\emph{A remark on quasi-isometries},
Proc. Amer. Math. Soc. 131 (2003), 1225--1231.
\end{thebibliography}

\end{document}
